\chapter{Gonorrhea Testing}

\section*{What Is This Test?}
Gonorrhea testing detects \textit{Neisseria gonorrhoeae}, which can cause urethritis, cervicitis, PID, proctitis, pharyngitis, conjunctivitis, and disseminated infection. Many infections are asymptomatic. Testing must cover the anatomic sites exposed.

\section*{Alternative Names}
GC test; gonorrhea NAAT; GC PCR; gonococcal culture; GC culture and susceptibility.

\section*{Principle}
NAAT is the preferred diagnostic method for validated genital and extragenital specimens. Culture remains essential when treatment failure or antimicrobial resistance is suspected because it permits antimicrobial susceptibility testing. Gram stain can rapidly support diagnosis in symptomatic men with urethral discharge but is not suitable for screening or for pharyngeal and rectal infection \cite{workowski2021sexually}.

\section*{History}
Culture on selective media was historically central to diagnosis. Molecular assays now provide more sensitive testing of urine and swab specimens, while culture remains important for resistance surveillance and suspected treatment failure.

\section*{Why This Test Is Ordered}
\begin{itemize}
  \item Screening according to age, pregnancy, sexual practices, risk, and local recommendations.
  \item Evaluation of urethritis, cervicitis, PID, proctitis, pharyngitis after exposure, conjunctivitis, or suspected disseminated infection.
  \item Testing of exposed partners and evaluation after a known diagnosis.
  \item Culture and susceptibility testing when treatment failure or resistance is suspected.
\end{itemize}

\section*{Primary vs Secondary Test}
NAAT is the primary screening and diagnostic test at validated sites. Culture is primary when susceptibility testing is required. Expedited partner therapy, where permitted, must follow current public-health guidance and the partner's clinical circumstances; it is not a universal automatic cefixime prescription.

\section*{How This Test Is Performed}
Collect first-catch urine or a vaginal, endocervical, pharyngeal, rectal, conjunctival, or other specimen according to anatomy and exposure. Use the manufacturer's transport medium for NAAT. Collect culture specimens into appropriate media and transport promptly because gonococci are fragile. Label every specimen with the anatomic site.

\section*{What Preparation Is Needed}
Fasting is not required. Avoid urinating for approximately 1 hour before first-catch urine when the laboratory recommends this. Recent antibiotic exposure should be documented but is not a general reason to defer diagnostic testing. If test of cure is indicated, use the timing and specimen method in current guidance.

\section*{Contraindications and Precautions}
There are no absolute contraindications to urine or swab collection. A negative result applies only to the tested site. Routine test of cure is not recommended for uncomplicated urogenital or rectal gonorrhea treated with a recommended regimen. Pharyngeal gonorrhea requires test of cure at 7--14 days after treatment. Suspected treatment failure requires culture, preferably with simultaneous NAAT, and antimicrobial susceptibility testing. Do not interpret a positive NAAT alone as evidence of viable organisms immediately after treatment.

\section*{Risks}
Urine collection is noninvasive. Swab collection may cause brief discomfort or spotting. Untreated infection can cause PID, infertility, ectopic pregnancy, disseminated infection, neonatal disease, and increased HIV transmission risk.

\section*{What Happens After the Test}
Treatment, partner management, public-health reporting, and abstinence instructions should follow current local guidance. Chlamydia treatment is added when chlamydia has not been excluded or is otherwise indicated. Test of cure is performed for pharyngeal infection at 7--14 days and for other situations specified by current guidance, such as suspected treatment failure or a nonstandard regimen.

\section*{Understanding Results}
\subsection*{Below Normal}
\begin{itemize}
  \item \textbf{Not detected:} No gonococcal nucleic acid was detected at the tested site; untested sites may still be infected.
  \item \textbf{Culture negative:} No organism was recovered; sensitivity is reduced by delayed transport, prior antibiotics, low organism burden, and pharyngeal infection.
\end{itemize}

\subsection*{Above Normal}
\begin{itemize}
  \item \textbf{NAAT detected:} Gonorrhea is detected at the tested site and requires treatment according to current guidance.
  \item \textbf{Culture positive:} Confirms infection and permits susceptibility testing.
  \item \textbf{Persistent symptoms or positive test of cure:} Obtain culture and susceptibility testing and consult infectious-disease/public-health experts.
\end{itemize}

\section*{Normal Results}
\textit{N. gonorrhoeae} not detected at each site tested and no organism isolated by culture.

\section*{Follow-up Tests}
\begin{itemize}
  \item Pharyngeal test of cure at 7--14 days.
  \item Culture and antimicrobial susceptibility testing for suspected treatment failure or resistance.
  \item Chlamydia, HIV, syphilis, and other STI testing according to exposure.
  \item Retesting for reinfection approximately 3 months after treatment.
\end{itemize}

\section*{References}
\bibliographystyle{unsrtnat}
\bibliography{references}
\clearpage
\section*{Regulatory Status and Authorization Date}
Regulatory status applies to the specific assay, instrument, manufacturer, and intended use rather than to the test name alone. No single approval date can be assigned to this test category. For a commercial assay, verify the current product in the relevant regulator's database: the U.S. Food and Drug Administration (FDA) clearance, approval, or authorization; India's Central Drugs Standard Control Organisation (CDSCO) licence or permission; the European Union In Vitro Diagnostic Regulation (EU IVDR) CE certificate issued through the applicable conformity-assessment route; or the United Kingdom Medicines and Healthcare products Regulatory Agency (MHRA) registration and UKCA/CE status. Laboratory-developed tests may not have an individual FDA, CDSCO, EU IVDR, or MHRA approval date.
\clearpage
